\documentclass[12pt]{article}

\usepackage{fullpage}
\usepackage{graphicx}
\usepackage{amssymb}
\usepackage{amsmath}
\usepackage[none]{hyphenat}
\usepackage{parskip}
\usepackage[spanish]{babel}
\usepackage[utf8]{inputenc}
\usepackage{hyperref}
\usepackage{fancyhdr}
\usepackage{tasks}
\usepackage{mdframed}
\usepackage{xcolor}
\usepackage{pgfplots}
\usepackage[makeroom]{cancel}
\usepackage{multicol}
\usepackage[shortlabels]{enumitem}
\usepackage{stackrel}
\usepackage{tkz-tab}
\usepackage{xpatch}
\usepackage{tkz-euclide}
\usetkzobj{all}
\xpatchcmd{\tkzTabLine}{$0$}{$\bullet$}{}{}

\setlength{\headheight}{10pt}
\setlength{\headsep}{10pt}
\pagestyle{fancy}
\rhead{\ayudantia \ - \alumno}
\tikzset{t style/.style={style=solid}}

\newcommand*{\mybox}[2]{\colorbox{#1!30}{\parbox{.98\linewidth}{#2}}}

\newenvironment{solucion}
{\begin{mdframed}[backgroundcolor=black!10]
		{\bf Solución:}\\
	}
	{
	\end{mdframed}
}

\newenvironment{alternativas}[1]
{\begin{multicols}{#1}
		\begin{enumerate}[a)]
		}
		{
		\end{enumerate}
	\end{multicols}
}

\newenvironment{preguntas}
{\begin{enumerate}\itemsep12pt
	}
	{
	\end{enumerate}
}

\newcommand{\ayudantia}{{\sc Ayudantía 15}}
\newcommand{\tituloayu}{Compilado Álgebra Lineal}
\newcommand{\fecha}{22 de junio de 2019}
\newcommand{\sigla}{MAT1203}
\newcommand{\nombre}{Álgebra Lineal}
\newcommand{\profesor}{Camilo Perez}
\newcommand{\ano}{2019}
\newcommand{\semestre}{1}
\newcommand{\mail}{mat1203@ifcastaneda.cl}
\newcommand{\alumno}{Ignacio Castañeda - \mail}

\newcommand{\ev}{\Big|}
\newcommand{\ra}{\rightarrow}
\newcommand{\lra}{\leftrightarrow}
\newcommand{\N}{\mathbb{N}}
\newcommand{\R}{\mathbb{R}}
\newcommand{\Exp}[1]{\mathcal{E}_{#1}}
\newcommand{\List}[1]{\mathcal{L}_{#1}}
\newcommand{\EN}{\Exp{\N}}
\newcommand{\LN}{\List{\N}}
\newcommand{\comment}[1]{}
\newcommand{\lb}{\\~\\}
\newcommand{\eop}{_{\square}}
\newcommand{\hsig}{\hat{\sigma}}
\newcommand{\widesim}[2][1.5]{
	\mathrel{\overset{#2}{\scalebox{#1}[1]{$\sim$}}}
}
\newcommand{\wsim}{\widesim{}}
\newcommand{\lh}{\stackrel{L'H}{=}}

\begin{document}
\thispagestyle{empty}

\begin{minipage}{2cm}
	\includegraphics[width=2cm]{../../../../img/logo.pdf}
	\vspace{0.5cm}
\end{minipage}
\begin{minipage}{\linewidth}
	\begin{tabular}{lrl}
		{\scriptsize\sc Pontificia Universidad Catolica de Chile} & \hspace*{0.7in}Curso: &
		\sigla  - \nombre\\
		{\sc Facultad de Matemáticas}&
		Profesor: & \profesor \\
		{\sc Semestre \ano-\semestre} & Ayudante: & {Ignacio Castañeda}\\
		& {Mail:} & \texttt{\mail}
	\end{tabular}
\end{minipage}

\vspace{-10mm}
\begin{center}
	{\LARGE\bf \ayudantia}\\
	\vspace{0.1cm}
	{\tituloayu}\\
	\vspace{0.1cm}
	\fecha\\
	\vspace{0.4cm}
\end{center}

\begin{preguntas}
\item Sea $A$ y $B$ dos matrices de $m\times n$ y $n \times p$, respectivamente. Demuestre que si las columnas de $B$ son linealmente dependientes, entonces también lo son las columnas de $AB$.
\item Sea $T: \R^2 \ra \R^3$ la transformación lineal definida por
	$$ T\left(\begin{bmatrix}
	x_1\\
	x_2
	\end{bmatrix}\right) = \begin{bmatrix}
	5x_1 - 3x_2\\
	4x_1 -x_2\\
	2x_1 +3x_2
	\end{bmatrix}$$
	encuentre $x = \begin{bmatrix}
	x_1\\
	x_2
	\end{bmatrix}$ tal que $T(x) = \begin{bmatrix}
	13\\
	9\\
	1
	\end{bmatrix}$
\item Determine condiciones sobre $\alpha$ y $\beta$ de modo que \\
	$\begin{pmatrix}
	\alpha-\beta\\
	\alpha + \beta\\
	2\alpha - \beta
	\end{pmatrix}$ pertenezca a $Gen\left\{\begin{pmatrix}
	-1\\
	-1\\
	1
	\end{pmatrix}\, \begin{pmatrix}
	1\\
	2\\
	1
	\end{pmatrix}, \begin{pmatrix}
	1\\
	3\\
	3
	\end{pmatrix}\right\}$
\item Describa todas las soluciones del siguiente sistema de ecuaciones y comparelas con las de su sistema homogeneo
	$$
	\begin{array}{rcr}
	x_1 +2x_2-3x_3& = & 5\\
	2x_1 + x_2 - 3x_3& = & 13\\
	-x_1 + x_2 & = & -8
	\end{array}$$
\item En las siguientes matrices aumentadas, determinar los valores de $h$ para que el sistema sea consistente.
\begin{tasks}(2)
\task $
		\begin{bmatrix}
		1 & h & 4\\
		3 & 6 & 8
		\end{bmatrix}
		$
\task $
		\begin{bmatrix}
		-4 & 12 & h\\
		2 & -6 & -3
		\end{bmatrix}
		$
\end{tasks}
\item Realiza las siguientes operaciones con los vectores dados
	$$
	v_1 = \begin{pmatrix}
	1\\
	7\\
	8
\end{pmatrix};\qquad
	v_2 = \begin{pmatrix}
	-4\\
	2\\
	-3
\end{pmatrix}; \qquad
	v_3 = \begin{pmatrix}
	0\\
	1\\
	3
\end{pmatrix}; \qquad
	v_4 = \begin{pmatrix}
	13\\
	-3\\
	1
\end{pmatrix},
	 $$
\begin{tasks}(4)
\task $v_1 + v_2$
\task $v_3 - v_4$
\task $v_2 \cdot v_3$
\task $v_1 \times v_2$
\task $v_2 \times v_1$
\task $(v_1 \times v_2) \cdot v_2$
\task $3 - ((2v_1) \cdot v_2)$
\task $v_2 - v_3 \times v_1$
\end{tasks}
\item Verificar si los siguientes puntos son o no colineares entre si
\begin{tasks}(2)
\task $P(3,2,5), \ P(0, 1/2, -1), \ P(5, 3, 9)$
\task $P(1,1,4), \ P(-2,-3.-8), \ P(4,4,16)$
\end{tasks}
\item Sean los vectores
		$$
	v_1 = \begin{pmatrix}
	2\\
	1\\
	5
\end{pmatrix};\qquad
	v_2 = \begin{pmatrix}
	-1\\
	-3\\
	0
\end{pmatrix}$$
\begin{enumerate}[a)]
\item Buscar un vector $v_3$ que sea perpendicular a $v_1$ y $v_2$ y verificar que efectivamente sea perpendicular a ambos.
\item Buscar un vector paralelo a $(v_2-v_3)$ y verificar que lo sea
\item Encontrar un vector perpendicular a $(v_3 + 3v_1)$
\end{enumerate}
\item Detereminar si los siguientes vectores son paralelos, perpendiculares u oblicuos
\begin{tasks}(2)
\task $$ 
			\begin{pmatrix}
			-1\\
			2\\
			5
		\end{pmatrix};\qquad
			\begin{pmatrix}
			2\\
			-4\\
			-10
		\end{pmatrix}$$
\task  $$ 
			\begin{pmatrix}
			2\\
			5\\
			10
		\end{pmatrix};\qquad
			\begin{pmatrix}
			-3\\
			8\\
			-11
		\end{pmatrix}$$
\end{tasks}
\item Determinar el plano que pasa por los puntos $P_1(4,-1,-2)$, $P_2(0,0,1)$ y $P_3(2,-3,0)$.
\item Dado $P_1(0,2,-3)$ y $P_2(1,0,3)$, determinar la recta que pasa por $P_1$ y $P_2$.
\item Encontrar las ecuaciones de dos planos diferentes cuya intersección sea la recta que pasa por el punto $P_1(1,3,-2)$ y $P_2(2,0,4)$.
\item Encontrar un plano que sea perpendicular al plano cuya ecuación es $3x -7y +2z = 5$ y que pase por el punto $P(0,2,-1)$
\item Determinar la ecuación de un plano que pase por $P(4,2,1)$ y que sea paralelo al plano de ecuación $2x-5y+z=6$
\item 
\begin{enumerate}[a)]
\item Sean $a=\left(\begin{array}{r}
      1\\-1\\1
    \end{array}\right)$ y $b=\left(\begin{array}{r}
      0\\-1\\2
    \end{array}\right)$, calcule el vector\textbf{ $b-proy_a b$ }y verifique que este es ortogonal a $a$.
\item Demuestre que si $a$ y $b\in\R^3$ el vector $b-proy_a b$ es ortogonal a $a$.
\end{enumerate}
\item Escribe los siguientes sistemas de ecuaciones en notación matricial.
\begin{tasks}(2)
\task $
		\begin{array}{rcr}
		x_1 -x_2 + 3x_3 & = & 2\\
		3x_2 - 2x_3 & = & 0\\
		4x_1 +2x_2 & = & -1
		\end{array}
		$
\task $
		\begin{array}{rcr}
		3x_1 & = & 1\\
		x_1-x_2& = & 3\\
		-2x_1+x_2 & = & -8\\
		x_2 & = & 0
		\end{array}
		$
\task $
		\begin{array}{rcr}
		2x_2 -x_3 +2x_4& = & 5\\
		x_1+x_2 - x_3 & = & 2\\
		-4x_1 +x_3-x_4& = & -4
		\end{array}
		$
\task $
		\begin{array}{rcr}
		x_1 +x_2-3x_3& = & 0\\
		x_2+4x_4& = & -1\\
		x_1-2x_2+x_3-5x_4 & = & 2\\
		x_3-x_4 & = & 5
		\end{array}
		$
\end{tasks}
\item Dadas las siguientes matrices aumentadas, escribir el sistema de ecuaciones asociado.
\begin{tasks}(2)
\task $
		\begin{bmatrix}
		-4 & 2 & 1 &  0\\
		2  & 1 & 2 & -1\\
		-3 & 5 &-6& -2
		\end{bmatrix}
		$
\task $
		\begin{bmatrix}
		2 & 1 & -3 & 0\\
		1  & -2 & 0 & 12\\
		4 & 2 &1& 8
		\end{bmatrix}
		$
\end{tasks}
\item Determinar si el siguiente sistema es consistente o no
	$$
	\begin{array}{rcr}
	x_1 -6x_2& = & 5\\
	x_2-4x_3+x_4& = & 0\\
	-x_1+6x_2+x_3+5x_4& = & 3\\
	-x_2+5x_3+4x_4 & = & 0
	\end{array}
	$$
\item Sea $A$ una matriz de $5 \times 3$ tal que existe una matriz $C$ de $3 \times 5$ tal que $CA = I_3$ y sea $b \in R^5$ tal que la ecuación $Ax = b$ tiene solución.\\
Demuestra que $Ax = b$ tiene solución única.
\item Lleve las siguientes matrices ampliadas a su forma escalonada reducida y determine la existencia y unicidad de las soluciones del sistema.
\begin{tasks}(2)
\task $
		\begin{bmatrix}
		1 & 2 & 4 & 8\\
		2 & 4 & 6 & 8\\
		3 & 6 & 9 & 12
		\end{bmatrix}
		$
\task $
		\begin{bmatrix}
		1 & 2 & 4 & 5\\
		2 & 4 & 5& 4\\
		4 & 5 & 4 & 2
		\end{bmatrix}
		$
\end{tasks}
\item Determine si el siguiente sistema de ecuaciones posee solución única, infinitas soluciones o no tiene solución
	$$
	\begin{array}{rcr}
	x+y-z& = & 1\\
	3x+2y+z& = & 1\\
	5x+3y+4z& = & 2\\
	-2x -y +5z & = & 6
	\end{array}$$
\item Sea la matriz $A=
	\begin{bmatrix}
	3 & 5 & -4\\
	-3 & -2 & 4\\
	6 & 1 & -8
	\end{bmatrix}$
\begin{enumerate}[a)]
\item Determinar el conjunto solución de su sistema homogeneo.
\item Describir todas las soluciones de $Ax=b$ con $b=
		\begin{pmatrix}
		7\\
		-1\\
		-4
		\end{pmatrix}$ 
\end{enumerate}
\item Determine todas las matrices $A$ tales que
	$$ A\begin{bmatrix}
	3\\
	2
	\end{bmatrix}=\begin{bmatrix}
	2\\
	1
	\end{bmatrix}$$
\item Considere el siguiente sistema de ecuaciones, donde $a$ es una constante.
  $$\begin{array}{llll}
   x_1&+x_2&+x_3&=1 \\
    x_1&+x_2&+ax_3&=1 \\
     ax_1&+ax_2&+x_3&=a\\
      x_1&-ax_2&+ax_3&=0  
  \end{array}$$
\begin{enumerate}[a)]
\item Determine valores de $a$ para los cuales el sistema es inconsistente.
\item Determine valores de $a$ para los cuales el sistema es consistente, y encuentre la solución.
\end{enumerate}
\item Sea la matriz $A=
	\begin{bmatrix}
	1 & 3 & 4\\
	-4 & 2 & -6\\
	-3 & -2 & -7
	\end{bmatrix}
	$ y $b$ un vector en $R^3$. ¿La ecuación $Ax=b$ es consistente para todo $b$?
\item Determinar las condiciones en $a$ para que la matriz
	$$A = \begin{bmatrix}
	a & 2a & 0 & 0 \\
	0 & 1 & 0 & 3a-1\\
	0 & 1 & a-1 & 2a-1\\
	a & 2a & 0 & a
	\end{bmatrix}$$
	sea invertible
\item Sean los vectores 
	$$
	v_1 = \begin{pmatrix}
	1\\
	1\\
	2
	\end{pmatrix};\qquad
	v_2 = \begin{pmatrix}
	2\\
	2\\
	-3
	\end{pmatrix}; \qquad
	v_3 = \begin{pmatrix}
	0\\
	1\\
	3
	\end{pmatrix}$$
	determinar si $Gen\{v_1, v_2, v_3\} = R^3$.
\item Sean 
$$v_1=\left(\begin{array}{r}
  1\\0\\-2
\end{array}\right), \quad v_2=\left(\begin{array}{r}
  -2\\1\\7
\end{array}\right) \quad y \quad v_3=\left(\begin{array}{r}
  h\\0\\-2
\end{array}\right)$$
¿Para qué valor(es) de $h$ $Gen\{v_1, v_2, v_3\}=Gen\{v_1,v_2\}$ ?
\item Sean $\{u, v, w\}$ un conjunto de vectores linealmente independientes. Demuestre que el conjunto $\{u+v, u+2w, v+3u+w\}$ es linealmente independiente.
\item Demuestre que el conjunto $\{u, v, w\}$ es L.I. si y solo si el conjunto $\{u+v, u+w, v+w\}$ es L.I.
\item Sea $A$ una matriz de $4 \times 4$ tal que $A=\begin{bmatrix}a_1 & 2a_1 & a_2 & a_1-a_2\end{bmatrix}$ con $a_1, a_2 \in \R^4$ vectores linealmente independientes, y sea $B$ una matriz inyectiva de $4 \times 2$ tal que $A \cdot B = 0$. Determine una matriz con las características de $B$.
\item Sea $A$ una matriz de $n \times n$ tal que existe una matriz $B \neq 0$ tal que $AB = 0$.\\
Demuestre que $A$ no es invertible.
\item Sea $T: \R^7 \ra \R^9$ una transformación lineal. Suponga que el conjunto $\{u,v\}$ es linealmente independiente en $\R^7$ pero $\{T(u), T(v)\}$ es linealmente dependiente. Demuestre que la ecuación $T(x)=0$ admite soluciones no triviales.
\item Sea $L: P_2(\R) \ra \R^2$ una transformación lineal tal que
	$$L(1+x) = \begin{pmatrix}
	1\\
	1
	\end{pmatrix}, \quad L(1+x+x^2) = \begin{pmatrix}
	1\\
	-1
	\end{pmatrix}\ y \ L(1+2x) = \begin{pmatrix}
	1\\
	2
	\end{pmatrix}$$
	determine $L(a+bx+cx^2)$ para todo $a,b,c \in \R$.
\item Sea $T$ una transformación lineal tal que
	$$T\begin{bmatrix}2 \\ 1\end{bmatrix} = \begin{bmatrix}1 \\ 0 \\ 1 \\ 0\end{bmatrix} 
	\quad y \quad
	T \begin{bmatrix}3 \\ 1\end{bmatrix}  =\begin{bmatrix}0 \\ 1 \\ 0 \\ 1\end{bmatrix}$$
	Determine la matriz que representa a $T$
\item Sea la transformación lineal $T$, dada por la matriz
	$$T = \begin{bmatrix}
	1 & 2 & -1\\
	2 & 4 & -2
	\end{bmatrix}$$
	Determine la imagen del plano de ecuación $x_1 + x_3 = 1$ por T.
\item Sea $\mathbb{P} \subset \R^3$ el plano de ecuación $x_1-x_2+1=0$ y sea $T:\R^3 \ra \R^4$ la transformación dada por $T(x) = Ax$, donde
$$A = \begin{bmatrix}
1 & 0 & 3\\
-1 & 1 & 0 \\
0 & 1 & 3\\
0 & -2 & -6
\end{bmatrix}$$
Demuestre que la imagen del plano $\mathbb{P}$ bajo la transformación $T$ es una recta en $\R^4$.
\item Determine si las siguientes afirmaciones son verdaderas o falsas.
\begin{enumerate}[a)]
\item Para todo $u, v \in \R^n$ se tiene que $||u+v|| < ||u||+||v||$
\item Si $u \in Gen\{v_1, v_2\}$ entonces $Gen\{u, v_1\} \subset Gen\{u, v_2\}$
\item Sea $A$ una matriz de $n \times m$ tal que $n < m $, entonces las columnas de A son L.D.
\item Si $A$ es una matriz de $2 \times 3$ tal que $Ax = \begin{pmatrix}
		1\\
		1
		\end{pmatrix}$ y $ Ax = \begin{pmatrix}
		1\\
		2
		\end{pmatrix}$ tienen solución, entonces las filas de $A$ son L.I.
\end{enumerate}
\item Sea
	$$A = \begin{bmatrix}
	1 & 0 & 6\\
	-1 & 2 & 0\\
	0 & 5 & -1
	\end{bmatrix}$$
	Escriba $A$ como un producto de matrices elementales
\item Determinar la matriz inversa de
$$A = \begin{bmatrix}
1 & -1 & 0 \\
0 & 1 & 0 \\
2 & 0 & 1
\end{bmatrix}$$
\item Sea $A$ una matriz de $n\times m$ con columnas LI. Demuestre que si $P=A(A^tA)^{-1}A^{t}$, entonces $P^2=P$ y $P=P^t$.
\item Sea $A$ una matriz de $3\times 3$ tal que
      $$A\left(\begin{array}{r}
  1\\0\\1
\end{array}\right)= \left(\begin{array}{r}
  1\\0\\0
\end{array}\right),\quad\quad A \left(\begin{array}{r}
  1\\-2\\4
\end{array}\right)= \left(\begin{array}{r}
  1\\0\\1
\end{array}\right), \qquad  A\left(\begin{array}{r}
  -1\\1\\1
\end{array}\right)= \left(\begin{array}{r}
  0\\1\\1
\end{array}\right).$$ Calcule $A^{-1}$.
\item Demuestre que $T(x_1,x_2, x_3)=(x_1,x_1+x_2, x_1+x_2+x_3)$ es invertible y encuentre una f\'ormula para $T^{-1}$.
\item Determine si las siguientes afirmaciones son verdaderas o falsas. En caso de ser verdaderas demuéstrelas y si son falsas de un contraejemplo.
\begin{enumerate}[a)]
\item Si $\{v_1, v_2, v_3\}$ es un conjunto linealmente dependiente de vectores en $\R^3$ entonces el conjunto $\{v_1, v_2\}$ también es linealmente dependiente.
\item Si $A$ es una matriz de $3 \times 3$ entonces la imagen del plano $x+y+z=0$ bajo la transformación $T(x)=Ax$ es un plano.
\item Si $A_{2\times 2} = (a_{ij}), \quad B_{2\times 3} = (b_{ij}), \quad AB = C = (c_{ij})$ tal que
$$a_{ij} = (-2)^{i+j}, \quad b_{ij} = (-3)^{i-j}$$
entonces $c_{23} = -\dfrac{56}{9}$
\end{enumerate}
\item En cada caso, determine si la afirmación es verdadera o falsa y justifique su respuesta.
\begin{enumerate}[a)]
\item Si el conjunto $\{v_1, v_2, v_3, v_4\}$ es L.I., entonces $\{v_1, v_2, v_3\}$ también lo es.
\item Si $A$ es una matriz de $3\times 5$ y $T$ es la transformación lineal definida por $T(x) = Ax$, entonces $T$ tiene dominio $\R^3$.
\item Un sistema con más ecuaciones que incógnitas es siempre consistente.
\end{enumerate}
\item Encuentre una factorziación $A = LU$ de
	$$A = \begin{bmatrix}
	2 & 4 & -1 & 5 & -2 \\
	-4 & -5 & 3 & -8 & 1 \\
	2 & -5 & -4 & 1 & 8 \\
	-6 & 0 & 7 & -3 & 1
	\end{bmatrix}$$
\item Sea $A = \begin{bmatrix}
	1 & 2 & 3\\
	2 & 8 & 4 \\
	3 & 4 & 11\end{bmatrix}$. Calcule la factorización $A=LDL^T$ de la matriz $A$ y en base a esto encuentre una matriz $R$ tal que $A = RR^T$.
\item Sea una descomposción $PA = LU$ donde
	$$P = \begin{bmatrix}
	0 & 0 & 1\\
	0 & 1 & 0\\
	1 & 0 & 0
	\end{bmatrix}, \quad
	L = \begin{bmatrix}
	1 & 0 & 0\\
	-1 & 1 & 0\\
	0 & 1 & 1
	\end{bmatrix}, \quad
	U = \begin{bmatrix}
	2 & 1 & 0\\
	0 & 1 & 1\\
	0 & 0 & -1
	\end{bmatrix}$$
	determine usando directamente esta factorización la tercera fila de $A^{-1}$.
\item Sea $A \in M_4 (\R)$ tal que $|A|=5$, encuentre
\begin{tasks}(4)
\task $|2^3A|$
\task $|A^5|$
\task $|-A|$
\task $||A|A^{-1}|$
\end{tasks}
\item Sea $A$ una matriz de $4\times 4$, tal que $det(A) = \alpha \neq 0$
\begin{enumerate}[a)]
\item Calcule  $det(5A)+det(3A^{-1})$ en terminos de $\alpha$.
\item Se sabe que $det(Adj(A))=8$. Calcule $\alpha$.
\end{enumerate}
\item Sea $P$ el paralelepípedo con un vertice en el origen y vertices adyacentes en $(1,4,0)$, $(-2,-5,2)$, $(-1,2,-1)$
\begin{enumerate}[a)]
\item Determinar el volumen de $P$.
\item Se define $T: \R^3 \ra \R^3$ como la transformación lineal definida por $$T(x,y,z)=(x+y,y+z,x-z)$$
		Calcule el volumen de $T(P)$.
\end{enumerate}
\item Sea $A$ un matriz tal que
$$A^3 = 
\begin{bmatrix}
2 & 0 & 2\\
-1 & 1 & 2\\
2 & -1 & -1
\end{bmatrix}$$
\begin{enumerate}[a)]
\item Calcule el determinante de la matriz $A^3$ mediante desarrollo por cofactores.
\item Demuestre que $A$ no es invertible.
\end{enumerate}
\item Sea $A$ una matriz de $n \times n$ definida por
	$$a_{i,j} = 
	\begin{cases}
	3 \quad si\ i\leq j \\
	1 \quad si\ i > j
	\end{cases}$$
	Calcule $det(A)$.
\item Calcular el determinante de siguiente matriz de $n\times n$
	$$ \begin{bmatrix}
	-3 & x & x & \dots & x \\
	x & -3 & x & \dots & x \\
	x & x & -3 & \dots & x \\
	\vdots & \vdots & \vdots & \ddots & \vdots \\
	x & x & x & \dots & -3
	\end{bmatrix}$$
\item Sean $v_1,v_2,v_3,v_4 \in \R^4$, encuentre el $|A|$ si\\
$A=[v_1-v_3+v_4 \;\;\;\;-v_2-v_3 \;\;\;\; v_3-v_1
\;\;\;\;v_1+v_2+2v_4 ]\in M_4 (\R)$
\item Determine la segunda columna de la inversa de $A = \left[ \begin{array}{rrr} 2&1&0\\ 0&1&2\\ 1&2&0 \end{array} \right]$ utilizando el metodo de cramer.
\item Sea $V = \mathbb{P}_2$ el espacio vectorial de los polinomios con coeficientes reales de grado menor o igual a 2. Demuestre que el conjunto
	$$U=\{p(x) \in V : p(1) + p(0) = p(-1)\}$$
	Demuestre que $U$ es subespacio de $V$.
\item Sea $V = \mathbb{P}_3$ y sea
$$W = \{p(x) \in V: p(1) = p(0) =0\}$$
Demostrar que $W$ es subespacio de $V$.
\item La matriz $A = \begin{bmatrix}
	1 & 3 & 4 & -1 & 2\\
	2 & 6 & 6 & 0 & -3\\
	3 & 9 & 3 & 6 & -3\\
	3 & 9 & 0 & 9 & 0
	\end{bmatrix}$ es equivalente por filas a $B = \begin{bmatrix}
	1 & 3 & 4 & -1 & 2\\
	0 & 0 & 1 & -1 & 1\\
	0 & 0 & 0 & 0 & -8\\
	0 & 0 & 0 & 0 & 0
	\end{bmatrix}$. Sin calcular bases, determine las dimensiones de Nul $A$, Col $A$, Col $A^T$ y Nul $A^T$.
\item En cada caso, determine si la afirmación es verdadera o falsa y justifique su respuesta.
\begin{enumerate}[a)]
\item El subconjunto de $\R^2$
		$$E=\{(x_1, x_2) \in \R^2:9x_1^2+2x_2^2 \leq 4\}$$
		es un subespacio vectorial de $\R^2$.
\item Si $A$ es una matriz de $n \times n$ tal que $A^2=A$, entonces $Col(A) \cap Nul(A) = \{0\}$.
\item Si $n$ es impar y $A$ es una matriz de $n\times n$ que satisface $A^T = -A$ entonces $Nul(A) \neq \{0\}$.
\end{enumerate}
\item Sea 
	$$U=\{p(x) \in \mathbb{P}_2 : p(1) + p(0) = p(-1)\}$$
	Determinar una base de $U$.
\item Sea $A$ una matriz tal que $dim(Nul(A)) = 3$ y
	$$A^T \begin{pmatrix}
	1 \\ -2 \\ 1 \\ 1
	\end{pmatrix} = \begin{pmatrix}
	0 \\ 0 \\ 0 \\ 0 \\ 0
	\end{pmatrix}, \quad A^T \begin{pmatrix}
	2 \\ -1 \\ -1 \\ 1
	\end{pmatrix} = \begin{pmatrix}
	0 \\ 0 \\ 0 \\ 0 \\ 0
	\end{pmatrix}$$
\begin{enumerate}[a)]
\item Determine las dimensiones de $Col(A)$, de $Fila(A)$ y de $Nul(A^T)$
\item Determine, en los casos que sea posible, bases para los espacios $Nul(A^T)$ y $Fila(A)$.
\end{enumerate}
\item Sean $B = \begin{bmatrix} 1 & -1 \\ 0 & 0 \end{bmatrix}$ y $C = \{X \in M_2(\R)|BX=0\}$. Determine una base y la dimensión de $C$.
\item Use vectores de coordenadas para probar la independencia lineal del conjunto de polinomios $\{1-2t^2-t^3,t+2t^3,1+t-2t^2\}$.
\item Sean $p_1(t) = 1-t^2, p_2(t) = 1+t, p_3(t) = 1+t+t^2$. Se sabe que $\{p_1(t), p_2(t), p_3(t)\}$ es una base de $\mathbb{P}_2$.
\begin{enumerate}[a)]
\item Exprese los polinomios $f(t) = 3-5t + 2t^2$ y $g(t) = 1-3t$ como combinaciones lineales de $\{p_1(t), p_2(t), p_3(t)\}$.
\item Use los vectores de coordenadas encontrados en la parte anterior para determinar si el conjunto $\{p_1(t), f(t), g(t)\}$ es L.I. o L.D.
\end{enumerate}
\item Sea
$$T\left(\left[ \begin{array}{c}
a\\ b\\ c \end{array} \right] \right) = (a+b+c) + (a-b+2c)x +
(3b-c)x^2
$$
Determine una base para $Nul(T) $ y para $Im (T) $ y las dimensiones de estos subespacios.
\item Sea $T: \mathbb{P}_2 \ra \mathbb{P}_2$ una transformación lineal cuya matriz de transformación respecto a la base $B = \{1, t, t^2\}$ está dada por
$$[T]_B = \begin{bmatrix}
1 & -1 & 1\\
0 & 2 & -1\\
-1 & 3 & -2
\end{bmatrix}$$
Determine $T(t^2+1)$ y el espacio Nulo de $T$.
\item Sean $B_1 = \{1+x,1-x,1+x^2\}$ y $B_2$ bases de $P_2(\R)$ tales que para todo $p \in P_2(\R)$
	$$ \begin{bmatrix}
	1 & 0 & -1\\
	1 & 1 & 0 \\
	0 & 1 & 2
	\end{bmatrix} [p]_{B_1} = [p]_{B_2}$$
	Determine los polinomios que forman la base $B_2$.
\item Sean $B$ y $C$ bases de un espacio vectorial $V$ y $P = \begin{bmatrix}4 &-1 \\ 6 & -1\end{bmatrix}$ la matriz de cambio de coordenadas tal que $[v]_C = P[v]_B\ \forall v \in V$
\begin{enumerate}[a)]
\item Demuestre que el conjunto $W = \{v \in V: [v]_C = 2[v]_B\}$ es un subespacio de $V$.
\item Si $B=\{v_1,v_2\}$ determine una base para $W$ en términos de la base de $B$.
\end{enumerate}
\item Calcular los valores y vectores propios de
	$$ A = \begin{bmatrix} 1 & 2 & -1\\ 1 & 0 & 1\\ 4 & -4 & 5\end{bmatrix}$$
\item La matriz $\begin{bmatrix} 1 & -2 \\ 1 & 3 \end{bmatrix}$ actúa sobre $\mathbb{C}^2$. Determine los valores propios y una base para cada espacio propio de $\mathbb{C}^2$.
\item Sea $A$ una matriz invertible, y sea $\lambda$ un valor propio de $A$. Demuestre que $\lambda \neq 0$ y que $\dfrac{1}{\lambda}$ es un valor propio de $A^{-1}$.
\item Sea $A$ una matriz de $5 \times 5$ tal que $A^t = -A$ y $q$ el polinomio dado por $q(x) = 2-x^2+4x^3$. Demuestre que $2$ es valor propio de la matriz $q(A)$.
\item Sea $A$ una matriz de $2\times 2$ de rango 1 tal que $A\begin{bmatrix} 1 \\ 2\end{bmatrix} = \begin{bmatrix} 2 \\ 4 \end{bmatrix}$. ¿Es $A$ diagonalizable? Justifique.
\item Diagonalice la matriz
	$$M = \begin{bmatrix}
	1 & 0 & 0\\
	1 & 1 & 2 \\
	1 & 0 & 3
	\end{bmatrix}$$
	y encuentre una matriz $N$ tal que $N^3 = M$.
\item Diagonalice $A =
	\begin{bmatrix} 
	-1 & -2 & 2 \\
	0 & -1 & 0 \\
	0 & -2 & 1
	\end{bmatrix}$ y diagonalice $B = A^{10} + A - I$ 
\item Diagonalice ortogonalmente
	$$M = \begin{bmatrix}
	1 & 0 & 1 \\
	0 & 2 & 0 \\
	1 & 0 & 1
	\end{bmatrix}$$
\item Sean $u, v$ dos vectores ortogonales en $\R^n$ tales que $||u|| = 1, ||v|| = \sqrt[]{3/2}$. Demuestre que el conjunto 
	$$B = \{u-v, 3u+2v\}$$
	es ortogonal y encuentre las coordenadas del vector $4u - 9v$ respecto al conjunto $B$.
\item Demuestre que si $P$ es una matriz ortogonal de  $n \times n$, entonces para todo $x, y \in \R^n$ se tiene que $Px \cdot Py = x \cdot y$
\item Determine si las siguientes afirmaciones son Verdaderas o Falsas.
\begin{enumerate}[a)]
\item El espacio fila de $AB$ es subespacio del espacio fila de $B$.
\item Dada una base en un espacio vectorial $V$, el vector coordenado de un vector de $V$ con respecto a esa base, es único.
\item Si $A$ y $B$ son similares y $\lambda$ es valor propio de $A$ entonces $\lambda$ es también valor propio de $B$.
\end{enumerate}
\item Sea 
	$U = Gen\left\{
	\begin{bmatrix}1\\0\\1\\1\end{bmatrix},
	\begin{bmatrix}0\\1\\2\\1\end{bmatrix},
	\begin{bmatrix}1\\1\\3\\2\end{bmatrix}
	\right\}$.
\begin{enumerate}[a)]
\item Encuentre una base ortonormal de $U$.
\item Encuentre la distancia de $b = 
		\begin{bmatrix}1\\1\\1\\1\end{bmatrix}$ a $U$.
\end{enumerate}
\item Un cierto experimento genera los datos $(1,3)$, $(2,5)$ y $(3,4)$. Describa el modelo que da un ajuste de mínimos cuadrados de esos puntos mediante una recta de la forma $y = \beta_0 +  \beta_1x$
\item Sea $A$ la matriz simétrica que representa a la forma cuadrática
	$$Q(x) = 9x_1^2 + 7x_2^2 + 11x_3^2 - 8x_1x_2 + 8x_1x_3$$
\begin{enumerate}[a)]
\item Encuentre una matriz ortogonal $P$ tal que el cambio de variable $x = Py$ transforma la forma $x^TAx$ en una forma cuadrática sin productos cruzados.
\item Escriba la forma cuadrática en las nuevas variables y clasifíquela.
\end{enumerate}
\item Sea $A = \begin{bmatrix}1 & -3 \\ -3 & 9\end{bmatrix}$
\begin{enumerate}[a)]
\item Encuentre una matriz $L$ cuadrada, triangular inferior con numeros 1 en la diagonal, y una matriz diagonal $D$ tal que $A = LDL^T$.
\item Encuentre la segunda descomposición de Cholesky de $A$
\item Realice un cambio de variable $x=Py$ que transforme la forma cuadrática
		$$Q\left(\begin{bmatrix}x_1\\x_2\end{bmatrix}\right) = \begin{bmatrix}x_1& x_2\end{bmatrix}A\begin{bmatrix}x_1\\ x_2\end{bmatrix}$$
		en una sin términos con producto cruzado.
\end{enumerate}
\item Calcule la descomposición en valores singulares de la matriz $A = \begin{bmatrix}4 & 11 & 14\\ 8 & 7 & -2\end{bmatrix}$
\item Sea la matriz
	$A = \begin{bmatrix}
	1 & 0 \\1 & 1 \\
	-1 & 1
	\end{bmatrix}$
	, determine su descomposición en valores singulares.
\end{preguntas}
\end{document}